% Ampliación sustantiva de los comparadores determinantes BSD.
% Módulo destinado al capítulo BSD de la Parte IV.

\section{Comparators of the common determinantal unit}
\label{sec:amp-bsd-comparadores}

The root equality \eqref{eq:bsd-suc-raiz} fixes a single unit in the Kato
and Poitou--Tate channels. From that unit, the comparators
transporting the derived regulator to its scalar publication
are now constructed. No comparator chooses a normalization after learning the
leading term: all arise from morphisms of complexes, exact
sequences and already oriented bases.

\subsection{The chain of determinant lines}
\label{subsec:amp-bsd-lineas}

For a based perfect complex \(C\), the convention used is
\begin{equation}\label{eq:amp-bsd-linea-determinante}
 \lambda(C):=\bigotimes_i
 \det H^i(C)^{(-1)^i}.
\end{equation}
Let \(\mathcal C_E^K\), \(\mathcal C_E^{\mathrm{Iw}}\),
\(\mathcal C_E^{\mathrm{PT}}\) and \(\mathcal C_E^{\mathrm{loc}}\) be,
respectively, the complex published by the Kato channel, its based
Iwasawa model, the self-dual reduced Poitou--Tate complex and the finite--singular
localization cone. Their lines are linked by
\begin{equation}\label{eq:amp-bsd-cadena-lineas}
 \mathcal L_E^K
 \xrightarrow{\ \mathfrak c_{K/\mathrm{FERS}}\ }
 \mathcal L_E^{\mathrm{Iw}}
 \xrightarrow{\ \mathfrak c_{\mathrm{PT}}\ }
 \mathcal L_E^{\mathrm{ht}}
 \xrightarrow{\ \mathfrak c_{\mathrm{STS}}\ }
 \mathcal L_E^{\mathrm{fin}}
 \xrightarrow{\ \mathfrak c_{\mathrm{tors}}\ }
 \mathcal L_E^{\mathrm{ar}}
 \xrightarrow{\ \mathfrak c_{\infty}\ }
 \mathcal L_E.
\end{equation}
In this formula, the abbreviation \(\mathrm{STS}\) denotes the joint step of
Smith, Tamagawa and Tate--Shafarevich. All the arrows of
\eqref{eq:amp-bsd-cadena-lineas} are isomorphisms of oriented lines.

\subsection{Kato--FERS chain comparison from the common unit}
\label{subsec:amp-bsd-comparacion-cadenas-kato-fers}

The equivalence giving rise to the first comparator can be constructed before
taking determinants. Let \(F^q\mathcal C\) be the filtration by nonadic
depth, identified with the \(T\)-adic filtration through the carry, and
write
\begin{equation}\label{eq:amp-bsd-complejos-kato-iwasawa}
 \mathcal C_E^K
 :=\operatorname{Tot}\!\left(
 P_E^K\widehat\Pi_0G_{m,n}\right),
 \qquad
 \mathcal C_E^{\mathrm{Iw}}
 :=[V_\Lambda\xrightarrow{S_E(T)}W_\Lambda].
\end{equation}
Let
\(\operatorname{car}_T:G_{m,n}\to\mathcal C_E^{\mathrm{Iw}}\) be the carry
reader assigning to a route \(\gamma\) the generator determined by its
initial and final faces, multiplied by
\(T^{\operatorname{Carr}(\gamma)}\). Multiplication of concatenations in
the Iwasawa model is denoted by \(\mu_{\mathrm{Iw}}\). The Alexander--Whitney
diagonal and the publication \(P_E^K\) then define, before
taking cohomology or determinants, the chain map
\begin{equation}\label{eq:amp-bsd-phi-kato-fers}
\begin{aligned}
 \Phi_{K/\mathrm{FERS}}&:
 \mathcal C_E^K\longrightarrow\mathcal C_E^{\mathrm{Iw}},\\
 \Phi_{K/\mathrm{FERS}}
 &:=
 \mu_{\mathrm{Iw}}\circ
 \bigl(P_E^K\widehat\otimes\operatorname{car}_T\bigr)
 \circ\Delta_{\mathrm{AW}},\\
 \Phi_{K/\mathrm{FERS}}(z_K)&=z_{\mathrm{Iw}}.
\end{aligned}
\end{equation}
where \(z_{\mathrm{Iw}}\) is the root generator induced by \(1_E\). Thus,
each normalized simplex is sent to the sum of its Alexander--Whitney
cuts, preserving in each summand the two faces and the carry
power. The formula does not consult the analytic leading term.

\begin{lemma}[Generator-by-generator criterion for equivalence]
\label{lem:amp-bsd-bases-emparejadas-kato-fers}
In each degree \(i\), let \(\mathcal B_i^K\) be the finite basis of normalized
simplices, ordered by depth, faces and carry, and let
\(\mathcal B_i^{\mathrm{Iw}}\) be the basis formed by the generators with the
same initial and final faces in the Iwasawa model. For a candidate
map
\[
 \upsilon_i:\mathcal B_i^K\longrightarrow
 \mathcal B_i^{\mathrm{Iw}},
\]
define the degree-zero residual
\begin{equation}\label{eq:amp-bsd-residual-generador-graduado}
 \varepsilon_i(b)
 :=\Phi_{K/\mathrm{FERS}}^i(b)\bmod T-\upsilon_i(b).
\end{equation}
If \(\upsilon_i\) is a bijection and
\(\varepsilon_i(b)=0\) for every \(b\in\mathcal B_i^K\), then the two
completed modules of degree \(i\) are free of the same rank
\(d_i=|\mathcal B_i^K|\) and, for each generator,
\begin{equation}\label{eq:amp-bsd-phi-en-cada-generador}
 \Phi_{K/\mathrm{FERS}}^i(b)
 =\upsilon_i(b)
  +T\sum_{c\in\mathcal B_i^{\mathrm{Iw}}}
    \lambda_{c,b}(T)c,
 \qquad \lambda_{c,b}(T)\in\Lambda.
\end{equation}
In particular,
\(\operatorname{gr}^{F}\Phi_{K/\mathrm{FERS}}^i=I\) on all
generators, not only on the root unit.
\end{lemma}

\begin{proof}
The vanishing of \eqref{eq:amp-bsd-residual-generador-graduado} says that the
column of \(b\) has constant term \(\upsilon_i(b)\). The remaining
coefficients belong to \(T\Lambda\), giving
\eqref{eq:amp-bsd-phi-en-cada-generador}. Bijectivity identifies the
bases and fixes the same rank. In those bases, the reduction modulo \(T\) of the
matrix of \(\Phi^i\) is the identity; hence
\(\operatorname{gr}^{F}\Phi^i=I\).
\end{proof}

The root equality
\(\Phi_{K/\mathrm{FERS}}(z_K)=z_{\mathrm{Iw}}\) fixes the common normalization.
The expression of \(P_E^K\) on each normalized simplex, the bijection
\(\upsilon_i\) and the residuals \(\varepsilon_i(b)\) provide a
generator-by-generator check of the equivalence already constructed; they neither choose nor
govern the BSD conclusion.

\begin{lemma}[Filtered equivalence and normalization of the root]
\label{lem:amp-bsd-equivalencia-kato-fers}
For the chain map constructed in
\eqref{eq:amp-bsd-phi-kato-fers}, the map \(\upsilon_i\) is the
basis correspondence induced by the faces and carry of each
simplex. The Alexander--Whitney formula verifies in each degree the
criterion of lemma \ref{lem:amp-bsd-bases-emparejadas-kato-fers}; this is a
check of the equivalence already constructed, not an additional premise.
The map \(\Phi_{K/\mathrm{FERS}}\) is a based equivalence of
perfect complexes. In the bases ordered by depth, its matrices in
each degree have the form
\begin{equation}\label{eq:amp-bsd-matriz-kato-fers}
 [\Phi_{K/\mathrm{FERS}}^i](T)=I+TB_i(T),
 \qquad B_i(T)\in M_{d_i}(\Lambda).
\end{equation}
Therefore, the unit
\begin{equation}\label{eq:amp-bsd-rho-determinantal}
 \rho_{E,\ell}(T)
 :=\prod_i\det\!\left(I+TB_i(T)\right)^{(-1)^i}
 \in\Lambda^\times
\end{equation}
satisface \(\rho_{E,\ell}(0)=1\).
\end{lemma}

\begin{proof}
The simplicial identity
\(\partial\Delta_{\mathrm{AW}}
=(\partial\otimes I+(-1)^{\deg}I\otimes\partial)
\Delta_{\mathrm{AW}}\), together with the compatibility of
\(P_E^K\) with faces and degeneracies, proves that
\eqref{eq:amp-bsd-phi-kato-fers} commutes with the differentials. On the associated
graded object, lemma
\ref{lem:amp-bsd-bases-emparejadas-kato-fers} identifies all the
generators and gives the identity. In those bases,
\eqref{eq:amp-bsd-phi-en-cada-generador} is precisely
\eqref{eq:amp-bsd-matriz-kato-fers}.

The completeness of \(\Lambda=K[[T]]\) makes the series converge
\begin{equation}\label{eq:amp-bsd-inversa-kato-fers}
 (I+TB_i(T))^{-1}
 =\sum_{n\geq0}(-TB_i(T))^n;
\end{equation}
these inverses also commute with the differentials, so they form the
based inverse of \(\Phi_{K/\mathrm{FERS}}\). Taking the alternating
determinant yields \eqref{eq:amp-bsd-rho-determinantal}; evaluating at
\(T=0\) gives one. Thus, the normalization is not postulated from the leading
term: it is the constant coordinate of an equivalence that is the identity
on the root unit.
\end{proof}

The first comparator is the determinant of
\(\Phi_{K/\mathrm{FERS}}\). If \(\mathfrak z_K\) denotes the element
coming from the unit \(1_E\), its image is characterized by
\begin{equation}\label{eq:amp-bsd-comparador-kato}
 \mathfrak c_{K/\mathrm{FERS}}(\mathfrak z_K)
 =\left[T^{-d}\rho_{E,\ell}(T)\det S_E(T)\right]_{T=0},
 \qquad
 \rho_{E,\ell}(0)=1.
\end{equation}
The unit of \(\rho_{E,\ell}\) is oriented and equals one at the root; it therefore does not
alter the order or the basis. Theorem
\ref{thm:bsd-suc-bockstein} transforms \eqref{eq:amp-bsd-comparador-kato} into
\begin{equation}\label{eq:amp-bsd-kato-bockstein}
 \mathfrak c_{K/\mathrm{FERS}}(\mathfrak z_K)
 =R_{\mathrm{Boc}}^{\mathrm{der}}(S_E)
 =\tau_{\mathrm{reg}}\otimes\bigotimes_{q\geq1}\tau_q.
\end{equation}

The Poitou--Tate comparator is constructed page by page. For
\(q\geq1\), the isomorphism \(j_q\) of
\eqref{eq:bsd-suc-jq} and the reduced Bockstein
\(\bar\beta_q\) induce
\begin{equation}\label{eq:amp-bsd-comparador-pt-q}
 \mathfrak c_{\mathrm{PT},q}(\tau_q)
 :=\det\!\left(j_q\circ\bar\beta_q\right)
 =\det h_E^{(q)}.
\end{equation}
The regular direction is transported through
\(\det\bar S_E(0)=\tau_{\mathrm{reg}}\). Knudsen--Mumford gluing and the
Koszul signs defined in \eqref{eq:bsd-suc-rboc} provide the unique
oriented isomorphism
\begin{equation}\label{eq:amp-bsd-comparador-pt}
 \mathfrak c_{\mathrm{PT}}
 \left(\tau_{\mathrm{reg}}\otimes\bigotimes_{q\geq1}\tau_q\right)
 =\tau_{\mathrm{reg}}^{\mathrm{ht}}
  \otimes\bigotimes_{q\geq1}\det h_E^{(q)}.
\end{equation}
On the first stable page, the form \(h_E^{(1)}\) is the
Poincaré--Néron--Tate form on the free Mordell--Weil lattice; its determinant is
\begin{equation}\label{eq:amp-bsd-regulador-nt}
 \det h_E^{(1)}=\operatorname{Reg}(E)
\end{equation}
in the basis induced by the root unit.

\begin{lemma}[Poitou--Tate unfolding of the arithmetic degree]
\label{lem:amp-bsd-grado-pt}
Let
\[
 \operatorname{Flag}_{\mathrm{Boc}}(E)
 :=\bigoplus_{q\geq1}V_q
\]
the finite flag of positive pages, each with the basis inherited from the
root unit. Orthogonal reduction and the isomorphisms \(j_q\) construct a
based isomorphism
\begin{equation}\label{eq:amp-bsd-psi-pt}
 \Psi_{\mathrm{PT}}:
 \bigl(E(\mathbb Q)/E(\mathbb Q)_{\mathrm{tors}}\bigr)\otimes K
 \xrightarrow{\ \sim\ }\operatorname{Flag}_{\mathrm{Boc}}(E).
\end{equation}
In particular,
\begin{equation}\label{eq:amp-bsd-rango-bandera}
 \operatorname{rank}E(\mathbb Q)
 =\sum_{q\geq1}\dim_KV_q
 =\operatorname{ord}_T\det S_E(T).
\end{equation}
\end{lemma}

\begin{proof}
The reduced complex is filtered by the powers of the augmentation ideal. In the
quotient at level \(q\), the regular direction has already been separated by
\(\bar S_E(0)\), and the only surviving block is \(V_q\). The equivalence
\(X^{\mathrm{orth}}\) identifies the opposite block with its dual, while
\(j_q\bar\beta_q\) provides the perfect pairing that lifts the
basis of the quotient to the preceding level. Starting at the last nonzero page and
repeating this lifting yields \(\Psi_{\mathrm{PT}}\).

In the bases ordered by depth, the matrix of the resulting map is
block triangular and all its diagonal blocks are identities: in the
first block because \(z_K\) and \(z_{\mathrm{PT}}\) have the same root, and
in the others because \(j_q\) and \(\bar\beta_q\) are based isomorphisms.
Therefore \(\Psi_{\mathrm{PT}}\) is invertible and loses no page.
Taking dimensions and using
\eqref{eq:bsd-suc-carry-orden} yields
\eqref{eq:amp-bsd-rango-bandera}. On taking determinants, the same
lifting is exactly the gluing of
\eqref{eq:amp-bsd-comparador-pt}; thus, equality of ranks and equality of
lines come from a single based map.
\end{proof}

\subsection{Localization triangle, Smith rank and finite factors}
\label{subsec:amp-bsd-smith-sha}

For each place \(v\), the finite local condition is a based subcomplex
\(\mathcal C_{E,v}^{\mathrm{fin}}\subset\mathcal C_{E,v}\), and one defines
\(\mathcal C_{E,v}^{\mathrm{sing}}\) by the triangle
\begin{equation}\label{eq:amp-bsd-triangulo-local}
 \mathcal C_{E,v}^{\mathrm{fin}}
 \xrightarrow{\ i_v\ }\mathcal C_{E,v}
 \longrightarrow\mathcal C_{E,v}^{\mathrm{sing}}
 \longrightarrow\mathcal C_{E,v}^{\mathrm{fin}}[1].
\end{equation}
If \(\operatorname{res}_v\) denotes the global restriction, the localization
map is, at chain level,
\begin{equation}\label{eq:amp-bsd-mapa-localizacion}
 \ell_E:
 \mathcal C_E^{\mathrm{glob,red}}\oplus
 \bigoplus_v\mathcal C_{E,v}^{\mathrm{fin}}
 \longrightarrow\bigoplus_v\mathcal C_{E,v},
 \qquad
 \ell_E\bigl(x,(y_v)_v\bigr)
 :=\bigl(\operatorname{res}_v x-i_vy_v\bigr)_v .
\end{equation}
The definition by cone of the reduced Selmer complex gives the triangle
\begin{equation}\label{eq:amp-bsd-triangulo-localizacion}
 \mathcal C_E^{\mathrm{red}}
 \longrightarrow
 \mathcal C_E^{\mathrm{glob,red}}\oplus
 \bigoplus_v\mathcal C_{E,v}^{\mathrm{fin}}
 \xrightarrow{\ \ell_E\ }
 \bigoplus_v\mathcal C_{E,v}
 \longrightarrow\mathcal C_E^{\mathrm{red}}[1].
\end{equation}
Therefore, local--global exactness comes from the cone of a map written in
\eqref{eq:amp-bsd-mapa-localizacion}; it is not subsequently added as a relation
between cardinalities.

In \eqref{eq:amp-bsd-triangulo-localizacion} the common root summand
is canceled, which is the same in both channels by
\eqref{eq:bsd-suc-raiz}. The resulting complement is denoted by
\begin{equation}\label{eq:amp-bsd-cono-local-reducido}
 \mathcal C_E^{\mathrm{loc,red}}
 :=q_{\mathrm{root}}\operatorname{Cone}(\ell_E)[-1].
\end{equation}
After separating the free Mordell--Weil lattice and its dual through
\(X^{\mathrm{orth}}\), a choice of the already oriented integral bases
represents this complex by
\begin{equation}\label{eq:amp-bsd-complejo-smith}
 \mathcal C_E^{\mathrm{loc,red}}
 \simeq[M_E\xrightarrow{\ D_E\ }N_E].
\end{equation}

\begin{lemma}[Rational acyclicity and full rank]
\label{lem:amp-bsd-aciclicidad-rango}
The complex of \eqref{eq:amp-bsd-complejo-smith} is acyclic after
tensoring with \(\mathbb Q\). In particular,
\[
 D_E\otimes\mathbb Q:M_E\otimes\mathbb Q
 \xrightarrow{\ \sim\ }N_E\otimes\mathbb Q
\]
is an isomorphism; \(M_E\) and \(N_E\) have the same rank and \(D_E\) has
full rank.
\end{lemma}

\begin{proof}
Let \(x\) be a cycle of the complex
\(\mathcal C_E^{\mathrm{loc,red}}\otimes\mathbb Q\). The reduction has
eliminated the summand \(p_{\mathrm{root}}\), hence
\(q_{\mathrm{root}}x=x\). In the oriented integral bases used in
\eqref{eq:amp-bsd-complejo-smith}, the finite--singular complement
decomposes into the normal blocks of
\cref{lem:bsd-suc-forma-normal}. In each block \(\lambda\), the root
coefficient is already zero and the cycle is written
\[
 x_{\lambda}=u_{\lambda}e_{\lambda}
                  +v_{\lambda}b_{\lambda},
 \qquad
 \partial x_{\lambda}
   =(u_{\lambda}-3c_{\lambda}v_{\lambda})g_{\lambda}=0.
\]
Since \(g_{\lambda}\) belongs to the free basis, we have
\(u_{\lambda}=3c_{\lambda}v_{\lambda}\), and the explicit formula for the
differential gives
\[
 x_{\lambda}
   =v_{\lambda}(3c_{\lambda}e_{\lambda}+b_{\lambda})
   =\partial(v_{\lambda}f_{\lambda}).
\]
The sum is finite in each degree, so
\(x=\partial\bigl(\sum_{\lambda}v_{\lambda}f_{\lambda}\bigr)\).
Every cycle is therefore a boundary by the same normal form that produced
the filler \eqref{eq:bsd-suc-hstar}, without introducing
\(H_{\mathrm{FS}}\) as a premise. The equality
\(p_{\mathrm{root}}z_K=z_{\mathrm{PT}}\) is what allows the root
direction to be canceled before this calculation; without it a degree-zero
class would remain. The self-dual equivalence
\(X^{\mathrm{orth}}:\mathcal C_E^{\mathrm{red}}\simeq
D_3(\mathcal C_E^{\mathrm{red}})\) pairs the two remaining degrees and
transports their bases with complementary opposite orientation. From this one
obtains \(\operatorname{rank}M_E=\operatorname{rank}N_E\).
Finally, acyclicity of a two-term complex of equal rank
is equivalent to \(D_E\otimes\mathbb Q\) being invertible.
\end{proof}

The lemma makes it possible to apply the Smith normal form without presupposing any of
its divisors. There exist oriented changes of basis \(U,V\) such that
\begin{equation}\label{eq:amp-bsd-morfismo-smith}
 U D_E V=\operatorname{diag}(s_1,\ldots,s_t),
 \qquad s_i\in\mathbb Z\setminus\{0\}.
\end{equation}
Therefore,
\begin{equation}\label{eq:amp-bsd-grupo-smith}
 F_E:=\operatorname{coker}D_E
 \simeq\bigoplus_{i=1}^{t}\mathbb Z/|s_i|\mathbb Z,
 \qquad
 |F_E|=\prod_{i=1}^{t}|s_i|<\infty.
\end{equation}

\begin{lemma}[Finite localization sequence]
\label{lem:amp-bsd-exactitud-sha-tamagawa}
The torsion fragment of the long cohomology sequence of
\eqref{eq:amp-bsd-triangulo-localizacion} is
\begin{equation}\label{eq:amp-bsd-exacta-sha-tamagawa}
 0\longrightarrow\operatorname{Sha}(E)
 \xrightarrow{\ \iota_{\mathrm{Sha}}\ }F_E
 \xrightarrow{\ \partial_{\mathrm{loc}}\ }
 \bigoplus_v\Phi_v(\mathbb F_v)
 \longrightarrow0,
 \qquad c_v=|\Phi_v(\mathbb F_v)|.
\end{equation}
\end{lemma}

\begin{proof}
A global cocycle represents a class of
\(\operatorname{Sha}(E)\) precisely when its restriction at each place
admits a primitive in
\(\mathcal C_{E,v}^{\mathrm{fin}}\). The class of the pair formed by that
cocycle and its local primitives in the cone of
\eqref{eq:amp-bsd-mapa-localizacion} defines
\(\iota_{\mathrm{Sha}}\). If this class is a boundary, the global component is a
boundary and the Tate--Shafarevich class is zero; thus,
\(\iota_{\mathrm{Sha}}\) is injective.

The map \(\partial_{\mathrm{loc}}\) takes a class of the cone, restricts it to
each quotient
\(\mathcal C_{E,v}^{\mathrm{sing}}\) of
\eqref{eq:amp-bsd-triangulo-local} and preserves its component class. The
identification
\[
 H^0(\mathcal C_{E,v}^{\mathrm{sing}})_{\mathrm{tors}}
 \simeq\Phi_v(\mathbb F_v)
\]
is induced by the finite--singular quotient, not by the order of the group.
A class of the cone has zero residue exactly when its local
representatives can be chosen in the finite subcomplexes; by the preceding
definition, that kernel is the image of \(\iota_{\mathrm{Sha}}\).
Finally, the last arrow of
\eqref{eq:amp-bsd-triangulo-local} lifts each component class to the
local complex. The exactness of the cone
\eqref{eq:amp-bsd-triangulo-localizacion}, after removing the two
free summands paired by \(X^{\mathrm{orth}}\), turns that
lifting into a preimage under \(\partial_{\mathrm{loc}}\). This proves
surjectivity and, consequently, all of
\eqref{eq:amp-bsd-exacta-sha-tamagawa}.
\end{proof}

Now, and only after lemma
\ref{lem:amp-bsd-aciclicidad-rango}, the group \(F_E\) is finite. The
injection of \eqref{eq:amp-bsd-exacta-sha-tamagawa} then proves the
finiteness of \(\operatorname{Sha}(E)\). Taking orders in that exact
sequence yields
\begin{equation}\label{eq:amp-bsd-cardinal-smith}
 |F_E|=|\operatorname{Sha}(E)|\prod_v c_v.
\end{equation}
This is the comparator
\(\mathfrak c_{\mathrm{STS}}\): it transports the Smith elementary factors
to the Tate--Shafarevich and Tamagawa factors while preserving their order in the line,
instead of inserting their cardinalities once the leading term has been calculated.

The torsion part is constructed by applying the determinant to the Mordell--Weil
lattice and its Pontryagin dual. The two finite factors are dual and
contribute
\begin{equation}\label{eq:amp-bsd-comparador-torsion}
 \mathfrak c_{\mathrm{tors}}:quad
 u\longmapsto
 \frac{u}{|E(\mathbb Q)_{\mathrm{tors}}|^2}.
\end{equation}
The Archimedean comparator is obtained by integrating the oriented Néron differential
over the based real cycle:
\begin{equation}\label{eq:amp-bsd-comparador-periodo}
 \Omega_E:=\int_{E(\mathbb R)}|\omega_E|,
 \qquad
 \mathfrak c_\infty:quad u\longmapsto\frac{u}{\Omega_E}.
\end{equation}
The differential and the cycle are fixed before evaluating \(L(E,s)\); therefore
\eqref{eq:amp-bsd-comparador-periodo} does not retrospectively normalize the
analytic section.

\subsection{Causal separation of primitives, identities and consequences}
\label{subsec:amp-bsd-separacion-causal}

The generator-by-generator criterion of
lemma \ref{lem:amp-bsd-bases-emparejadas-kato-fers} audits in each degree the
equivalence already constructed. It does not condition its existence. The logical order
of the package is
\begin{equation}\label{eq:amp-bsd-diagrama-causal}
\begin{array}{c}
 \begin{gathered}
  1_E,\ \Delta_{\mathrm{AW}},\ P_E^K,\ P_E^{\mathrm{PT}},
  \ h_*=-vf,\ X^{\mathrm{orth}},\\
  \{\,\mathcal C_{E,v}^{\mathrm{fin}}\to\mathcal C_{E,v}\,\}_v,\\
  \text{Mordell--Weil and Pontryagin bases;}\\
  \text{Néron differential and cycle}
 \end{gathered}
 \\[2mm]\Downarrow\\[-1mm]
 \begin{gathered}
  \Phi_{K/\mathrm{FERS}}\text{ equivalence},\quad
  \rho_{E,\ell}(0)=1,\quad
  p_{\mathrm{root}}z_K=z_{\mathrm{PT}},\\
  R_{\mathrm{Boc}}^{\mathrm{der}}=\operatorname{LT}_T(S_E),\quad
  \det(j_q\bar\beta_q)=\det h_E^{(q)},\\
  H^*(\mathcal C_E^{\mathrm{loc,red}}\otimes\mathbb Q)=0,\quad
  D_E\otimes\mathbb Q\text{ isomorphism},\\
  0\to\operatorname{Sha}(E)\to F_E
  \to\bigoplus_v\Phi_v(\mathbb F_v)\to0
 \end{gathered}
 \\[2mm]\Downarrow\\[-1mm]
 \begin{gathered}
  d_{\mathrm{an}}=\operatorname{rank}E(\mathbb Q),\qquad
  \operatorname{Reg}(E),\qquad
  |\operatorname{Sha}(E)|\prod_vc_v,\\
  |E(\mathbb Q)_{\mathrm{tors}}|^{-2},\qquad
  \Omega_E^{-1},\qquad
  \text{equality of the leading term}.
 \end{gathered}
\end{array}
\end{equation}
The first row contains only already based objects and maps, together with
the indicated generator-by-generator verification; the second, identities
proved by those maps; the third, their consequences on the determinant
line. In particular, no factor in the last row is
used to retroactively normalize an arrow in the first.

\begin{proposition}[Exact package of comparators]
\label{prop:amp-bsd-paquete-comparadores-exacto}
The comparators of
\eqref{eq:amp-bsd-cadena-lineas} come from a single package of maps of
complexes and satisfy, in the order of construction,
\begin{equation}\label{eq:amp-bsd-paquete-identidades}
 \begin{gathered}
 \Phi_{K/\mathrm{FERS}}(z_K)=z_{\mathrm{Iw}},\qquad
 \operatorname{gr}^{F}\Phi_{K/\mathrm{FERS}}=I,\\
 \det\nolimits_{\mathrm{alt}}\Phi_{K/\mathrm{FERS}}
   =\rho_{E,\ell}(T),\qquad \rho_{E,\ell}(0)=1,\\
 \det(j_q\bar\beta_q)=\det h_E^{(q)}\quad(q\geq1),\\
 \Psi_{\mathrm{PT}}:
  \bigl(E(\mathbb Q)/E(\mathbb Q)_{\mathrm{tors}}\bigr)\otimes K
  \xrightarrow{\sim}\operatorname{Flag}_{\mathrm{Boc}}(E),\\
 p_{\mathrm{root}}z_K=z_{\mathrm{PT}},\qquad
 h_*=-vf,\quad \partial h_*=z_{\mathrm{PT}}-z_K,\\
 H^*(\mathcal C_E^{\mathrm{loc,red}}\otimes\mathbb Q)=0,\qquad
 U D_E V=\operatorname{diag}(s_1,\ldots,s_t),\\
 0\longrightarrow\operatorname{Sha}(E)
  \longrightarrow F_E
  \longrightarrow\displaystyle\bigoplus_v\Phi_v(\mathbb F_v)
  \longrightarrow0 .
 \end{gathered}
\end{equation}
The induced composition on determinant lines is exactly
\begin{equation}\label{eq:amp-bsd-paquete-composicion}
 \det(\Phi_{K/\mathrm{FERS}})
 \xrightarrow{\ \det(j_q\bar\beta_q)\ }
 \det(D_E)
 \xrightarrow{\ \mathrm{Smith}\ }
 \frac{|\operatorname{Sha}(E)|\prod_vc_v}
      {|E(\mathbb Q)_{\mathrm{tors}}|^2\,\Omega_E},
\end{equation}
with the Knudsen--Mumford signs and the orientations of
\eqref{eq:amp-bsd-cadena-lineas}. In particular, the right-hand side of
\eqref{eq:amp-bsd-paquete-composicion} is the final coordinate of the transported
element, not data used to construct any of the arrows.
\end{proposition}

\begin{proof}
The first row of \eqref{eq:amp-bsd-paquete-identidades} is
\eqref{eq:amp-bsd-phi-kato-fers} and
\eqref{eq:amp-bsd-matriz-kato-fers}--
\eqref{eq:amp-bsd-rho-determinantal}. The second is
\eqref{eq:amp-bsd-comparador-pt-q} and lemma
\ref{lem:amp-bsd-grado-pt}. The third combines equality of the root with
the explicit filler \eqref{eq:bsd-suc-hstar}; its linear extension in the
normal form induces the auxiliary reader used in lemma
\ref{lem:amp-bsd-aciclicidad-rango}. On restricting to the complement of the
root, that reduction contracts every rational cycle. The last row is
\eqref{eq:amp-bsd-morfismo-smith} and
\eqref{eq:amp-bsd-exacta-sha-tamagawa}. Applying the determinant functor to
those identities, in that same order, gives
\eqref{eq:amp-bsd-paquete-composicion}. No equality is obtained
by first comparing the two final scalars.
\end{proof}

\begin{lemma}[Causal independence of the leading term]
\label{lem:amp-bsd-independencia-termino-principal}
Let \(\mathscr P_E\) be the data set
\begin{equation}\label{eq:amp-bsd-primitivas-paquete}
 \mathscr P_E:=
 \left(
  1_E,\Delta_{\mathrm{AW}},P_E^K,P_E^{\mathrm{PT}},
  h_*=-vf,X^{\mathrm{orth}},
  \{\mathcal C_{E,v}^{\mathrm{fin}}\xrightarrow{i_v}
     \mathcal C_{E,v}\}_v,
  \text{bases and orientations},
  \omega_E,\gamma_E^{\mathbb R}
 \right),
\end{equation}
where \(\omega_E\) is the Néron differential and
\(\gamma_E^{\mathbb R}\) the oriented real cycle. All maps of the
package in proposition
\ref{prop:amp-bsd-paquete-comparadores-exacto} are determined by
\(\mathscr P_E\). The following numerals are not part of \(\mathscr P_E\)
\[
 L^{(r)}(E,1),\quad r,\quad\operatorname{Reg}(E),\quad
 |\operatorname{Sha}(E)|,\quad(c_v)_v,\quad
 |E(\mathbb Q)_{\mathrm{tors}}|,\quad\Omega_E.
\]
Those values are obtained, respectively, by evaluating the analytic section,
taking the degree of the flag, calculating determinants and Smith forms,
taking orders of finite groups and integrating
\(\omega_E\) over \(\gamma_E^{\mathbb R}\).
\end{lemma}

\begin{proof}
The formula \eqref{eq:amp-bsd-phi-kato-fers} uses only
\(1_E,\Delta_{\mathrm{AW}},P_E^K\) and the carry. The Bockstein and
Poitou--Tate maps use the filtration, \(P_E^{\mathrm{PT}}\), the bases and
self-duality. The cone of \eqref{eq:amp-bsd-mapa-localizacion} uses the
maps \(i_v\) and the local restrictions; its contraction and Smith
form use the explicit filler \eqref{eq:bsd-suc-hstar}, its auxiliary linear
extension, \(X^{\mathrm{orth}}\) and the integral bases. The last
two comparators use the torsion lattices and the
pair \((\omega_E,\gamma_E^{\mathbb R})\). This is an exhaustive enumeration
of the arguments appearing in the construction formulas. The
numerals listed in the statement appear only after applying
cohomology, determinant, cardinality or integration, so they cannot
retrospectively select the maps.
\end{proof}

\subsection{Composition, degrees and leading term}
\label{subsec:amp-bsd-composicion}

Define the total comparator by composition in the written order:
\begin{equation}\label{eq:amp-bsd-comparador-total}
 \mathfrak C_E:=
 \mathfrak c_\infty\circ
 \mathfrak c_{\mathrm{tors}}\circ
 \mathfrak c_{\mathrm{STS}}\circ
 \mathfrak c_{\mathrm{PT}}\circ
 \mathfrak c_{K/\mathrm{FERS}}.
\end{equation}
Since \(p_{\mathrm{root}}z_K=z_{\mathrm{PT}}\), all arrows receive the
same based unit. The filler \(h_*=-vf\) of
\eqref{eq:bsd-suc-hstar} explicitly fills the difference; its linear
extension in the normal basis produces the auxiliary homotopic reader and proves that
no additional class modifies the element of the line during
transport. The reader is not a separate premise.

\begin{theorem}[Construction of the based comparators]
\label{thm:amp-bsd-comparadores-construidos}
For every elliptic curve \(E/\mathbb Q\) with the complexes, bases,
orientations and local data defined above, the maps
\eqref{eq:amp-bsd-phi-kato-fers},
\eqref{eq:amp-bsd-comparador-pt},
\eqref{eq:amp-bsd-triangulo-localizacion},
\eqref{eq:amp-bsd-comparador-torsion} and
\eqref{eq:amp-bsd-comparador-periodo}, together with lemmas
\ref{lem:amp-bsd-aciclicidad-rango} and
\ref{lem:amp-bsd-exactitud-sha-tamagawa}, construct the five comparators of
\eqref{eq:amp-bsd-cadena-lineas}. Their composition preserves the unit, degree
and orientation, and satisfies
\begin{equation}\label{eq:amp-bsd-igualdad-grados}
 d_{\mathrm{an}}=\operatorname{ord}_T\det S_E(T)
 =d_{\mathrm{ar}}=\operatorname{rank}E(\mathbb Q).
\end{equation}
Moreover, \(\operatorname{Sha}(E)\) is finite and the equality of the distinguished
element in \(\mathcal L_E\) is published as
\begin{equation}\label{eq:amp-bsd-termino-principal}
 \frac{L^{(r)}(E,1)}{r!\,\Omega_E}
 =\frac{|\operatorname{Sha}(E)|\,\operatorname{Reg}(E)
        \prod_v c_v}
       {|E(\mathbb Q)_{\mathrm{tors}}|^2},
 \qquad r=\operatorname{rank}E(\mathbb Q).
\end{equation}
\end{theorem}

\begin{proof}
The Kato--FERS comparator and the condition
\(\rho_{E,\ell}(0)=1\) identify the analytic order with
\(d=\operatorname{ord}_T\det S_E(T)\) and the initial term with
\(R_{\mathrm{Boc}}^{\mathrm{der}}(S_E)\). The identity
\eqref{eq:bsd-suc-lt} preserves the regular page and every positive page. The
isomorphisms \(j_q\) transport those pages to the derived heights.
Lemma \ref{lem:amp-bsd-grado-pt} constructs the isomorphism of the flag with the
free Mordell--Weil lattice and, together with
\eqref{eq:bsd-suc-carry-orden}, proves
\eqref{eq:amp-bsd-igualdad-grados}.

Lemma \ref{lem:amp-bsd-aciclicidad-rango} deduces rational acyclicity
directly from the explicit filler \eqref{eq:bsd-suc-hstar}, its linear
reduction of the complement, cancellation of the common root and
\(X^{\mathrm{orth}}\); from this the full rank of \(D_E\) is obtained, not
as an additional hypothesis. The Smith form
\eqref{eq:amp-bsd-morfismo-smith} then produces the finite group \(F_E\).
Lemma \ref{lem:amp-bsd-exactitud-sha-tamagawa} extracts from the triangles
\eqref{eq:amp-bsd-triangulo-local} and
\eqref{eq:amp-bsd-triangulo-localizacion} the finite exact sequence; only
afterwards are the finiteness of \(\operatorname{Sha}(E)\) and
\eqref{eq:amp-bsd-cardinal-smith} deduced. Finally,
\eqref{eq:amp-bsd-regulador-nt},
\eqref{eq:amp-bsd-comparador-torsion} and
\eqref{eq:amp-bsd-comparador-periodo} publish, respectively, the regulator,
the torsion quotient and the period. The root equality eliminates any relative
unit between the two publications. Evaluating the same based element at
both ends of \eqref{eq:amp-bsd-comparador-total} yields
\eqref{eq:amp-bsd-termino-principal}.
\end{proof}

\subsection{Premises and refutation criteria}
\label{subsec:amp-bsd-falsadores}

The construction starts from the genealogical unit \(1_E\) and the Alexander--Whitney
diagonal. It also uses the
coupled publications, all the
Bocksteins and the regular page, the explicit filler \(h_*=-vf\), the
reduced self-duality \(X^{\mathrm{orth}}\), the chain maps of the
localization triangle and the oriented bases of torsion and period. The
Kato--FERS equivalence and its normalization to one, the full rank of \(D_E\), the
exactness of \eqref{eq:amp-bsd-exacta-sha-tamagawa} and the finiteness of
\(\operatorname{Sha}(E)\) are conclusions of the preceding lemmas, not
independent premises. The generator-by-generator criterion of lemma
\ref{lem:amp-bsd-bases-emparejadas-kato-fers} is a nongoverning check
that locally audits the first comparator. The criteria that would refute a specific arrow
are:
\begin{enumerate}
 \item a root coefficient different from one, which would introduce a relative
       unit between \(z_K\) and \(z_{\mathrm{PT}}\);
 \item a noninvertible \(j_q\) or the omission of a Bockstein page, which
       would alter the degree or the height determinant;
 \item a zero Smith divisor, which would prevent the finiteness of \(F_E\);
 \item failure of exactness of
       \eqref{eq:amp-bsd-exacta-sha-tamagawa}, which would break the identification
       of the finite factors;
 \item an orientation reversal in torsion or period, which would change the
       based element even if it preserved its unoriented line.
\end{enumerate}
The memoryless Manin projection and visible rescaling modulo three
violate the first criterion; they therefore constitute necessity proofs for
those reduced models, not premises of theorem
\ref{thm:amp-bsd-comparadores-construidos}.
